\begin{pastexam}[label={px:ntu-109-dsaa-05}]{臺大資工所 109 學年度 第 5 題}
（複選；正確的選項可能有 0 到 5 個，都不對時答 none。 「previous question」是上一題第 4 題）

(3 points) We will follow the notations from the previous question. We can build a heap from an
array of $n$ keys by the following algorithm. A \textit{heapify} operation for a binary minimum
heap combines two heaps and a root into a new heap. For ease of discussion we assume that the
array is indexed from 1 to $n$, and the left/right child of a node with array index $i$ have
array indices $2i$ and $2i+1$ respectively. We first set each key of the second half of the
array (indexed from $n/2$ to $n$) as $n/2$ heaps, each has one node. Then we combine two of them
and their parent into a heap of three nodes, and so on. If the root is \textit{smaller} than the
roots of \textit{both} subtrees then we stop. Otherwise we exchange the root with the smaller
root from the subtrees, and repeatedly heapify the subtree where the root has just been
replaced. Next we heapify the key at $n/2-1$ and all the way back to the root (at index 1). The
final result will be a binary minimum heap of $n$ nodes.
\begin{enumerate}[label=\arabic*.]
\item The minimum number of key comparisons is $\Theta(n)$.
\item The maximum number of key comparisons is $\Theta(n)$.
\item Since each heapify on a key may compare it with all the keys down to a leaf, its number of
comparisons is no more than the height of the heap, which is $O(\log n)$. Also we have $O(n)$
keys to heapify, so the total number of key comparisons is $O(n\log n)$.
\item The number of key comparisons is in the same order as the sum of (\textit{original, before
heapified}) levels (as defined in the previous question, the distance to the root) of all keys.
\item The number of key comparisons is \textit{maximized} if and only if the \textit{largest}
$n/2$ keys are in the \textit{first} half of the array and in \textit{decreasing} order.
\end{enumerate}
\end{pastexam}

\begin{pastsoln}
答案：1、2、3。

\textbf{演算法和計數方式}。 這是由下而上建堆積（資結書第~8.2~節，虛擬碼~8.3），只是改成最小堆積。
原題說「後半段（$n/2$ 到 $n$）各是一個節點的堆積」、再從 $n/2 - 1$ 開始 heapify，差了一：精確地說，
葉節點是位置 $\lfloor n/2 \rfloor + 1$ 到 $n$，heapify 從 $\lfloor n/2 \rfloor$ 做到 $1$。 這不影響各選項
的判斷，下面照精確的版本。 heapify 一個節點時每往下一層：有兩個子節點就先比兩個子節點（$1$ 次），
再比根和較小的子節點（$1$ 次）；只有一個子節點就只比 $1$ 次。 根比較小就停，否則交換、往下一層。
所以每往下一層最多 $2$ 次比較。

令 $h(i)$ 是以位置 $i$ 為根的子樹的高度（用邊數算，葉節點是 $0$）。 heapify 位置 $i$ 最多往下
$h(i)$ 層，比較次數介於 $1$ 和 $2h(i)$ 之間。

\begin{itemize}
\item 1. 最少的比較次數：每一個非葉節點（共 $\lfloor n/2 \rfloor$ 個）heapify 時至少比 $1$ 次，所以至少
$\lfloor n/2 \rfloor$ 次；上界見選項 2，是 $O(n)$。 所以最少的次數是 $\Theta(n)$。 例如陣列本來就由小到大
排好時，每個節點比完子節點就停，$n = 2^k - 1$ 時剛好 $n - 1$ 次。 對。
\item 2. 最多的比較次數：總數 $\le \sum_i 2h(i)$。 資結書定理~8.1 證明了完整二元樹所有節點的高度和小於
$2n$（那裡的高度用節點數算，$h(x) - 1$ 就是這裡的 $h(x)$），所以最多不到 $4n$ 次，是 $O(n)$；又至少
$\lfloor n/2 \rfloor$ 次。 最大值是 $\Theta(n)$。 例如由大到小排好的陣列，每個節點都一路換到葉節點，
次數大約是 $2n$。 對。
\item 3. 這是建堆積的「粗略分析」：每個節點最多往下 $O(\log n)$ 層，一共 $O(n)$ 個節點，所以
$O(n \log n)$。 嚴格說，每往下一層有兩次比較（兩個子節點互比一次、和根比一次），一個 heapify 最多
比 $2\lfloor \lg n \rfloor$ 次，不是「不超過高度」；但只差常數倍，$O(\log n)$ 和最後的 $O(n \log n)$
都不受影響。 $O$ 只是上界，$O(n \log n)$ 雖然不緊（真正是 $\Theta(n)$，見選項 2），仍然是正確的敘述。
對。
\item 4. 原本各節點的 level 加總：位置 $i$ 的 level 是 $\lfloor \lg i \rfloor$，位置 $\lfloor n/2 \rfloor + 1$
到 $n$ 這一半的節點 level 都 $\ge \lfloor \lg n \rfloor - 1$，所以加總至少 $\frac{n}{2}(\lfloor \lg n
\rfloor - 1) = \Omega(n \log n)$；而比較次數是 $\Theta(n)$。 兩者不同階。 跟比較次數同階的是「各節點
\textbf{高度}的總和」（到葉節點的距離），不是深度（到根的距離）：深的節點多，但它們往下走不了幾層。 錯。
\item 5. 「若且唯若」兩個方向都不成立。 取 $n = 4$，heapify 位置 $2$（一個子節點，$1$ 次），再 heapify
位置 $1$（$2$ 次；換到位置 $2$ 時再 $1$ 次），所以最多 $4$ 次。 前半段是位置 $1$、$2$，「最大的兩個鍵
由大到小放在前半段」只有 $\langle 4, 3, 1, 2\rangle$ 和 $\langle 4, 3, 2, 1\rangle$ 兩種。
\begin{itemize}
\item $\langle 4, 3, 1, 2\rangle$：位置 $2$ 的 $3$ 和子節點 $2$ 比（$1$ 次），交換成
$\langle 4, 2, 1, 3\rangle$。 位置 $1$：$2$ 和 $1$ 比（$1$ 次），$4$ 和 $1$ 比（$1$ 次），換到位置 $3$，
位置 $3$ 是葉節點，停。 共 $3$ 次，不是最多。 所以「若」不成立。
\item $\langle 2, 1, 3, 4\rangle$：位置 $2$ 的 $1$ 和 $4$ 比（$1$ 次），不換。 位置 $1$：$1$ 和 $3$ 比、
$2$ 和 $1$ 比（$2$ 次），換到位置 $2$，再和 $4$ 比（$1$ 次）。 共 $4$ 次，是最多，但最大的兩個鍵
$3$、$4$ 在後半段。 所以「唯若」也不成立。
\end{itemize}
錯。 列舉 $n \le 8$ 的所有排列，達到最多次數的排列都很多（$n = 4$ 時 $24$ 個裡有 $12$ 個），遠不只
題目說的那一種。
\end{itemize}
\end{pastsoln}
